









It remains only to prove Lemma~\ref{lem:Gseq}.
\begin{proof}[Proof of Lemma~\ref{lem:Gseq}] We prove Lemma~\ref{lem:Gseq} by induction on $|P|$, with $P=\emptyset$ being obvious. 

Write $J(P)$ for the set of order ideals of $P$. 
We break up~\eqref{eq:signedcount} according to the first order ideal $I_1$ and proceed inductively on $P\setminus I_1$. Start with the equality
\begin{equation}\label{eq:first}
\sum_\text{$\mathcal{I}$ a $G$-sequence} \! (-1)^{|\mathcal{I}|} = \sum_{\emptyset\ne A \in J(P)} \,\,\sum_{\stackrel{\text{$\mathcal{I}$ a $G$-sequence}}{I_1=A}} (-1)^{|\mathcal{I}|}.
\end{equation}
Notice that for any order ideal $A$, the partition on the cover relations of $P\setminus A$ induced by $G\sqcup B$ is again acyclic. Then by induction, the nonzero contributions to the right hand side of~\eqref{eq:first} come from nonempty order ideals $A$ in which 
\begin{itemize}
\item all cover relations inside $A$ are good,
\item all cover relations inside $P\setminus A$ are bad.
\end{itemize}

Let $\mathcal{A}$ be the set of nonempty order ideals of $P$ satisfying these conditions. Then using~\eqref{eq:signedcount} inductively,~\eqref{eq:first} becomes
\begin{equation}\label{eq:second}
\sum_\text{$\mathcal{I}$ a $G$-sequence} \! (-1)^{|\mathcal{I}|} = \sum_{A\in \mathcal{A}} \,\,(-1)\cdot (-1)^{|P\setminus A|}.
\end{equation}

It remains to identify $\mathcal{A}$ in terms of $P$ and $G$, which we do as follows. Let $Y$ be the maximal up-closed subset of $P$ such that all cover relations within $Y$ are bad. Note that $Y$ is uniquely defined, since if $Y_1$ and $Y_2$ are up-closed subsets satisfying that condition, then $Y_1\cup Y_2$ also satisfies the condition. 

Let $X = P\setminus Y$. Let $Y'\subseteq Y$ be the subset consisting of the minimal elements $y\in Y$ satisfying that if $x\lessdot y$ then $(x,y)\in G$. 
Then we claim 
\begin{enonce}{claim}\label{c:A}\ \\*[-1.2em]
\begin{enumerate}
\item\label{claim3.20_1} If $X = \emptyset$ then $\mathcal{A} = \{I \subseteq \min(P): I\ne \emptyset\}$, where $\min(P)$ denotes the minimal elements of $P$.
\item\label{claim3.20_2} If $X\ne \emptyset$ and some cover relation within $X$ is in $B$, then $\mathcal{A}=\emptyset$.
\item\label{claim3.20_3} If $X\ne \emptyset$ and all cover relations within $X$ are in $G$, then
\[
\mathcal{A} = \{X\cup I: I\subseteq Y'\}.
\]
\end{enumerate}
\end{enonce}


\begin{proof}[Proof of Claim~\ref{c:A}.] If $X=\emptyset$ then $G=\emptyset$, and $\mathcal{A}$ then consists of all nonempty order ideals with no cover relations within them. So part~\eqref{claim3.20_1} follows.

Suppose $X\ne \emptyset$. Suppose $A\in \mathcal{A}$. Now for each maximal element $x\in X$, there is some $y\in Y$ such that $x\lessdot y$ is good. So necessarily $x\in A$, since otherwise the covering relation $x\lessdot y$ would lie in $P\setminus A$. So $A$ contains all maximal elements of $X$; thus $A\supseteq X$. So if some cover relation within $X$ is in $B$, then $\mathcal{A}=\emptyset$, proving part~\eqref{claim3.20_2}. 

Otherwise, we see that $X\in\mathcal{A}$. Furthermore, if $A\in\mathcal{A}$ then $A\cap Y$ must be an antichain in $Y$; otherwise $A$ contains a bad cover relation. So $A\cap Y\subseteq \min(Y)$. And if $y\in A\cap Y$, then any cover relation $x\lessdot y$ must be good. We conclude that $\mathcal{A} \subseteq \{X\cup I: I\subseteq Y'\}$; the reverse containment also clearly follows.
\end{proof}

Now from Claim~\ref{c:A}, the rest of the proof of Lemma~\ref{lem:Gseq} can be deduced from~\eqref{eq:second} by using the obvious identity 
$\sum_{S\subseteq T} (-1)^{|S|} = 0$
for nonempty finite sets $T$. Explicitly, in the case~\ref{c:A}\eqref{claim3.20_1}, Equation~\eqref{eq:second} becomes 
\[
\sum_{\emptyset\ne A \subseteq \min(P)} \,\,(-1)\cdot(-1)^{|P\setminus A|} = (-1)^{|P|}.
\]
In the case~\ref{c:A}\eqref{claim3.20_2}, Equation~\eqref{eq:second} is the empty sum. 
In the case~\ref{c:A}\eqref{claim3.20_3}, Equation~\eqref{eq:second} becomes
\[
\sum_{I\subseteq Y'} \,\,(-1)\cdot(-1)^{|P\setminus(X \cup I)|} = (-1)^{|P|+1}\sum_{I\subseteq Y'} (-1)^{|I|},
\]
which is $0$, as desired, provided that $Y'\ne \emptyset$. It remains only to show that $Y'=\emptyset$ is not possible.

If $Y'=\emptyset$ then by definition of $Y'$, every minimal element $y\in Y$ covers some $x\in X$ such that $x\lessdot y$ is bad. And every maximal element $x\in X$ is covered by some $y \in Y$ such that $x\lessdot y$ is good. Therefore in the orientation of the Hasse diagram of $P$ described in Definition~\ref{def:acyclic}, \emph{every} element in $X$ sends out an upwards edge, and \emph{every} element in $Y$ sends out a downwards edge. Therefore the oriented Hasse diagram has no sink and cannot be acyclic, contradicting the hypotheses.
\end{proof}



\section{An inverse formula}



\looseness-1
We give an analogous linear expansion of skew Schur functions into skew stable Grothendieck polynomials. This formula generalizes~\cite[Theorem~2.7]{lenart-combinatorial}, which pertains to straight shapes, and visibly specializes to that result when $\sigma$ is a straight shape. 

\begin{theorem}\label{t:back}
For any connected skew shape $\sigma$, 
\begin{equation}\label{eq:gsigma-conj}
s_\sigma = \sum_{\mu\supseteq \sigma} (-1)^{|A(\mu/\sigma)|} b_{\sigma,\mu} \cdot G_\mu
\end{equation}
where the $b_{\sigma, \mu}$ is the product of the following two numbers:
\begin{enumerate}
\item the number of row-weakly-bounded, reverse-row-strict tableaux of shape $A( \mu/\sigma)$, and 
\item the number of row-bounded, semistandard tableaux of shape $B(\mu/\sigma)$.
\end{enumerate}
\end{theorem}

\begin{proof}
Fix $\sigma$. By Theorem~\ref{thm:comb}, the right hand side of~\eqref{eq:gsigma-conj} is
\begin{equation}\label{eq:a-b}
\sum_{\mu\supseteq \sigma} \sum_{\lambda\supseteq \mu} (-1)^{|A(\mu/\sigma)|+|B(\lambda/\mu)|} a_{\lambda,\mu} \cdot b_{\mu,\sigma} \cdot s_\lambda
\end{equation}
and $a_{\lambda,\mu} \cdot b_{\mu,\sigma}$ is the product of the sizes of the following two sets:
\begin{enumerate}
\item\label{proof_theo4.1_1} the row-weakly-bounded tableaux $T'$ on $A(\lambda/\sigma)$ which are semistandard on $A(\lambda/\mu)$ and reverse-row-strict on $A(\mu/\sigma)$, and
\item\label{proof_theo4.1_2} the row-bounded tableaux $T''$ on $B(\lambda/\sigma)$ which are reverse-row-strict on $B(\lambda/\mu)$ and semistandard on $B(\mu/\sigma)$.
\end{enumerate} 

Now for a fixed $\lambda$, a row-weakly-bounded filling $T'$ of $A(\lambda/\sigma)$, and a row-bounded filling $T''$ of $B(\lambda/\sigma)$, the contribution of the triple $(\lambda,T',T'')$ to the sum~\eqref{eq:a-b} is
\begin{equation}\label{eq:should-vanish}
\sum_{\substack{\mu\text{ with } \sigma\subseteq \mu\subseteq \lambda \\T',T''\text{ satisfy }\eqref{proof_theo4.1_1},\eqref{proof_theo4.1_2}}}
 (-1)^{|A(\mu/\sigma)|+|B(\lambda/\mu)|}\cdot s_\lambda
 \end{equation}

Now if $\lambda = \sigma$ then~\eqref{eq:should-vanish} equals $s_\sigma$ vacuously. Otherwise, we show that~\eqref{eq:should-vanish} vanishes. In words, the coefficient of $s_\lambda$ in~\eqref{eq:should-vanish} is a signed count of ways to divide the two shapes $A(\lambda/\sigma)$ and $B(\lambda/\sigma)$, such that in each, the upper-left is semistandard and the lower-right is reverse-row-strict. Then the following lemma, applied to $A(\lambda/\sigma)$ and $B(\lambda/\sigma)$, finishes the proof of Theorem~\ref{t:back}. 
\end{proof}

\begin{lemma}\label{lem:ss-rrs}
Let $\mu$ be any nonempty skew shape. A \emph{division} of $\mu$ is a partition $\mu = \mu_S \sqcup \mu_R$ into two skew shapes such that $\mu_S$ is an order ideal of $\mu$ considered as a poset (Definition~\ref{def:young}). In other words, no box of $\mu_R$ is north/west of any box of $\mu_S$.
Suppose $T$ is a (non-set-valued) tableau on $\mu$. Say that $(\mu_S,\mu_R)$ is \emph{allowed} by $T$ if $T$ is semistandard on $\mu_S$ and reverse-row-strict on $\mu_R$. Then
\[
\sum_{(\mu_S,\mu_R) \text{ allowed by }T} (-1)^{|\mu_S|} = 0.
\]
\end{lemma}

\begin{exam}\label{ex:strip}
If $T = (a_1,\ldots,a_n)$ is a tableau on a single row of length $n$, then the above sum is empty unless $a_1 \le \cdots \le a_M > \cdots > a_n$ for a unique index $M$. In this case, there are two allowable divisions: where $\mu_S$ is either the first $M$ or the first $M-1$ boxes.
\end{exam}

\begin{proof}[Proof of Lemma~\ref{lem:ss-rrs}]
Each cover relation $b_1 \lessdot b_2$ in $\mu$ may be labelled $S$ or $R$ according to whether the restriction of $T$ to the $2$-box shape $\{b_1,b_2\}$ is semistandard or is reverse-row-strict. In the first case, write $S$ in box $b_1$; in the second case, write $R$ in box $b_2$. In this way, fill the boxes of $\mu$ using the alphabet $\{\emptyset, S, R, SR\}$. (In Example~\ref{ex:strip}, the first $M-1$ boxes are $S$, the next box is empty, and the remaining are $R$.) 

Now if there are any allowable divisions at all, then no box is labelled $SR$, the $S$ boxes must be an order ideal, the boxes containing $R$ must be the complement of an order ideal, and the empty boxes form an antichain in between. Then $(\mu_S,\mu_R)$ is allowable if and only if $\mu_S$ contains all $S$ boxes and $\mu_R$ contains all $R$ boxes. Then to prove the lemma, it is enough to show that there exists at least one empty box. Consider, among all boxes with maximum number in $T$, an upper-rightmost one. That box has empty label.
\end{proof}

\longthanks{We thank Vic Reiner, Bridget Tenner and Alexander Yong for generously answering some questions over email about their work~\cite{reiner-tenner-yong}. We also thank Jennifer Morse, Travis Scrimshaw, and an anonymous referee for helpful contextual remarks and references. We are grateful to Darij Grinberg for correcting an earlier error in the formulation of Lemma~\ref{lem:Gseq}, and for other detailed comments.}
